\usepackage[T1,T2A]{fontenc}
\usepackage[utf8]{inputenc}
\usepackage[english,main=russian]{babel}
\usepackage{fix-cm}

\usepackage[
	left=30mm,
	right=15mm,
	top=20mm,
	bottom=20mm,
]{geometry}


\usepackage{microtype}
\sloppy

\usepackage{setspace}
\onehalfspacing

\usepackage{indentfirst}
\setlength{\parindent}{12.5mm}

\makeatletter
\renewcommand\LARGE{\@setfontsize\LARGE{22pt}{20}}
\renewcommand\Large{\@setfontsize\Large{20pt}{20}}
\renewcommand\large{\@setfontsize\large{16pt}{20}}
\makeatother
\usepackage{titlesec}
\titleformat{\chapter}[block]{\hspace{\parindent}\large\bfseries}{\thechapter}{0.5em}{\large\bfseries\raggedright}
\titleformat{name=\chapter,numberless}[block]{}{}{0pt}{\large\bfseries\centering}
\titleformat{\section}[block]{\hspace{\parindent}\large\bfseries}{\thesection}{0.5em}{\large\bfseries\raggedright}
\titleformat{\subsection}[block]{\hspace{\parindent}\large\bfseries}{\thesubsection}{0.5em}{\large\bfseries\raggedright}
\titleformat{\subsubsection}[block]{\hspace{\parindent}\large\bfseries}{\thesubsection}{0.5em}{\large\bfseries\raggedright}
\titlespacing{\chapter}{12.5mm}{-22pt}{10pt}
\titlespacing{\section}{12.5mm}{10pt}{10pt}
\titlespacing{\subsection}{12.5mm}{10pt}{10pt}
\titlespacing{\subsubsection}{12.5mm}{10pt}{10pt}
\makeatletter
\renewcommand*{\l@chapter}[2]{
\ifnum \c@tocdepth>\m@ne
    \addpenalty{-\@highpenalty}
    \vskip 1em \@plus.2em
    \@dottedtocline{0}{0pt}{1.5em}{\bfseries #1}{\bfseries #2}
\fi
}

\makeatother

\usepackage{xcolor}
\usepackage{graphicx}
\usepackage{float}
\usepackage{wrapfig}
\usepackage{tikzscale}
\usepackage{pgfplots}
\pgfplotsset{compat=newest}
\newcommand{\includeimage}[5]{%
	\ifthenelse{\equal{#2}{f}}{%
		\begin{figure}[#3]
			\center{\includegraphics[width=#4]{inc/img/#1}}%
			\caption{#5}%
			\label{img:#1}%
		\end{figure}%
	}{%
		\ifthenelse{\equal{#2}{w}}{%
			\begin{wrapfigure}{#3}{#4}
				\center{\includegraphics[width=#4]{inc/img/#1}}%
				\caption{#5}%
				\label{img:#1}%
			\end{wrapfigure}%
		}{%
			\PackageError{bmstu}{unknown image type}{}%
		}%
	}%
}

\usepackage{listings}
\usepackage{listingsutf8}
\lstset{
	inputencoding=utf8/koi8-r,
	basicstyle=\small\ttfamily,
	rulecolor=\color{black},
	escapeinside={\%*}{*)},
	breaklines=true,
	breakatwhitespace=true,
	tabsize=4,
	showstringspaces=false,
	float=h!,
	abovecaptionskip=-5pt,
}
\definecolor{numbers}{rgb}{0.5,0.5,0.5}
\definecolor{keywords}{rgb}{0.13,0.13,1}
\definecolor{comments}{rgb}{0,0.5,0}
\definecolor{strings}{rgb}{0.9,0,0}
\newcommand{\includelisting}[2]{%
	\lstinputlisting[
		frame=single,
		caption={#2},
		label={lst:#1},
	]{inc/lst/#1}%
}
\newcommand{\includelistingpretty}[3]{%
	\lstinputlisting[
		language={#2},
		keywordstyle=\color{keywords},
		stringstyle=\color{strings},
		commentstyle=\color{comments},
		frame=leftline,
		numbers=left,
		numberstyle=\footnotesize\color{numbers},
		caption={#3},
		label={lst:#1},
	]{inc/lst/#1}%
}

\usepackage{tabularx}
\usepackage{booktabs}
\usepackage[
	labelsep=endash,
	figurename=Рисунок,
	singlelinecheck=false,
]{caption}
\captionsetup[figure]{justification=centering}

\usepackage{lscape}
\usepackage{afterpage}

\usepackage{amsmath}
\usepackage{amssymb}

\usepackage[
	style=gost-numeric,
	language=auto,
	autolang=other,
	sorting=none,
	movenames=false,
]{biblatex}
\usepackage{csquotes}
\DeclareFieldFormat{urldate}{(дата обращения:\addspace\thefield{urlday}\adddot \thefield{urlmonth}\adddot \thefield{urlyear})}

\usepackage[unicode,hidelinks]{hyperref}

\usepackage{xifthen}
\usepackage{enumitem}
\usepackage{lastpage}
\usepackage{totcount}
\usepackage{assoccnt}
\newcounter{appendixchapters}
\DeclareAssociatedCounters{chapter}{appendixchapters}
\regtotcounter{appendixchapters}

\makeatletter
\newcommand*{\bmstu@toc@dottedline}[5]{%
	\ifnum #1>\c@tocdepth \else
		\vskip \z@ \@plus.2\p@
		{%
			\leftskip #2\relax
			\rightskip \@tocrmarg plus 1fil\relax
			\parfillskip -\rightskip
			\parindent #2\relax
			\@afterindenttrue
			\interlinepenalty\@M
			\leavevmode
			\@tempdima #3\relax
			\advance\leftskip \@tempdima
			\null\nobreak\hskip -\leftskip
			{#4}\nobreak
			\leaders\hbox{$\m@th \mkern \@dotsep mu\hbox{.}\mkern \@dotsep mu$}\hfill
			\nobreak
			\hb@xt@\@pnumwidth{\hfil\normalfont\normalcolor #5}%
			\par
		}%
	\fi
}
\newcommand{\bmstuTocNoHyphenSetup}{%
	\hyphenpenalty=10000
	\exhyphenpenalty=10000
	\let\@dottedtocline\bmstu@toc@dottedline
}
\makeatother

\newcommand{\maketableofcontents}{%
	\renewcommand\contentsname{СОДЕРЖАНИЕ}%
	\begingroup
		\bmstuTocNoHyphenSetup
		\tableofcontents
	\endgroup
}
\newcommand{\makebibliography}{%
	\addcontentsline{toc}{chapter}{СПИСОК ИСПОЛЬЗОВАННЫХ ИСТОЧНИКОВ}%
	\printbibliography[title=СПИСОК ИСПОЛЬЗОВАННЫХ ИСТОЧНИКОВ]%
}

\newcommand{\makeappendix}[1]{%
	\chapter*{\hfill{\centering ПРИЛОЖЕНИЕ #1}\hfill}%
	\addcontentsline{toc}{chapter}{ПРИЛОЖЕНИЕ #1}%
	\renewcommand{\thelstlisting}{#1.\arabic{lstlisting}}%
	\setcounter{lstlisting}{0}%
}
\usepackage{lastpage}
\usepackage{etoolbox}
\usepackage{totcount}
\usepackage{assoccnt}
\usepackage{hyperref}

\newcounter{totfigures}
\newcounter{tottables}
\providecommand\totfig{}
\providecommand\tottab{}

\makeatletter
\AtEndDocument{%
	\addtocounter{totfigures}{\value{figure}}%
	\addtocounter{tottables}{\value{table}}%
	\immediate\write\@mainaux{%
		\string\gdef\string\totfig{\number\value{totfigures}}%
		\string\gdef\string\tottab{\number\value{tottables}}%
	}%
}
\makeatother

\pretocmd{\chapter}{%
	\addtocounter{totfigures}{\value{figure}}%
	\setcounter{figure}{0}%
}{}{}

\pretocmd{\chapter}{%
	\addtocounter{tottables}{\value{table}}%
	\setcounter{table}{0}%
}{}{}

\newcounter{totbibentries}
\newcommand*{\listcounted}{}
\providecommand\totbib{}

\makeatletter
\AtDataInput{%
	\xifinlist{\abx@field@entrykey}\listcounted
	{}
	{%
		\stepcounter{totbibentries}%
		\listxadd\listcounted{\abx@field@entrykey}%
	}%
}
\makeatother


\newenvironment{essay}[1]
{
	\chapter*{РЕФЕРАТ}
	\addcontentsline{toc}{chapter}{РЕФЕРАТ}

	Расчётно-пояснительная записка
	\begin{NoHyper}\pageref{LastPage}\end{NoHyper}~страниц,
	\totfig~рисунков,
	\tottab~таблиц,
	\thetotbibentries~источников,
	2~приложения

	\noindent \MakeUppercase{#1}\par
}{}

% --- definitions / abbreviations (из .sty) ---
\newcommand{\definition}[2]{%
	\item \noindent #1~---~#2
}

\newenvironment{definitions}
{
	\chapter*{ОПРЕДЕЛЕНИЯ}
	\addcontentsline{toc}{chapter}{ОПРЕДЕЛЕНИЯ}

	В настоящей расчётно-пояснительной записке применяют следующие термины с соответствующими определениями.

	\begin{description}[leftmargin=0pt]
	}
	{
	\end{description}
}

\newenvironment{abbreviations}
{
	\chapter*{ОБОЗНАЧЕНИЯ И СОКРАЩЕНИЯ}
	\addcontentsline{toc}{chapter}{ОБОЗНАЧЕНИЯ И СОКРАЩЕНИЯ}

	В настоящей расчётно-пояснительной записке применяют следующие сокращения и обозначения.

	\begin{description}[leftmargin=0pt]
	}
	{
	\end{description}
}

% \makeatletter
% \renewcommand{\@pnumwidth}{2.5em}   % было 1.55em по умолчанию
% \renewcommand{\@tocrmarg}{3.5em}    % правый отступ для оглавления
% \makeatother
% --- end ---
